% =====================================================================
%  FICHE 8 — THÈME 4 — NIVEAU FACILE — CORRIGÉ
% =====================================================================
\documentclass[11pt,a4paper]{article}
\usepackage[utf8]{inputenc}
\usepackage[T1]{fontenc}
\usepackage{lmodern}
\usepackage[french]{babel}
\usepackage{amsmath,amssymb,mathtools}
\usepackage{icomma}
\usepackage{siunitx}
\sisetup{output-decimal-marker={,},per-mode=power,group-separator={\,},group-minimum-digits=5}
\usepackage[version=4]{mhchem}
\usepackage{xcolor}
\usepackage{tikz}
\usepackage[most]{tcolorbox}
\usepackage{enumitem}
\usepackage{array}
\usepackage{fancyhdr}
\usepackage[hidelinks]{hyperref}
\usetikzlibrary{arrows.meta,positioning,calc}
\usepackage[a4paper,margin=2.1cm,headheight=26pt,headsep=10pt]{geometry}
\usepackage{titlesec}

\definecolor{primary}{RGB}{150,98,162}
\definecolor{primarydark}{RGB}{92,52,110}
\definecolor{excol}{RGB}{0,135,90}

\tcbset{boxbase/.style={enhanced,breakable,boxrule=0.9pt,arc=2.5pt,left=3mm,right=3mm,top=2.3mm,bottom=2.3mm,fonttitle=\bfseries,coltitle=white,toptitle=0.6mm,bottomtitle=0.6mm}}
\newtcolorbox[auto counter]{correction}[1][]{boxbase,colback=excol!4,colframe=excol,title={\textsc{Corrigé}~\thetcbcounter\ \ifx\relax#1\relax\relax\else\textmd{--- \itshape #1}\fi}}

\newcommand{\dS}{\Delta S}
\newcommand{\rep}[1]{\textcolor{primarydark}{\textbf{#1}}}

\pagestyle{fancy}
\fancyhf{}
\fancyhead[L]{\small\textcolor{primarydark}{\textbf{Fiche 8} — Thème 4 : Entropie et deuxième principe}}
\fancyhead[R]{\small\textcolor{primarydark}{\textsc{Facile — corrigé}}}
\fancyfoot[C]{\small\thepage}
\renewcommand{\headrulewidth}{0.8pt}
\renewcommand{\headrule}{\hbox to\headwidth{\color{primary}\leaders\hrule height \headrulewidth\hfill}}
\setlength{\parindent}{0pt}\setlength{\parskip}{3pt}

\begin{document}

\begin{center}
\begin{tikzpicture}
\node[rectangle,rounded corners=7pt,fill=excol,text=white,minimum width=16cm,
      minimum height=1.1cm,align=center,inner sep=6pt]
  {{\large\bfseries Thème 4 — Entropie et deuxième principe}\\[1pt]{\small Niveau \textsc{Facile} — corrigé détaillé}};
\end{tikzpicture}
\end{center}

\begin{correction}[ordre du désordre]
\begin{enumerate}[label=\alph*),leftmargin=2em,topsep=2pt,itemsep=2pt]
  \item Entropie croissante : \rep{glace $<$ eau liquide $<$ vapeur d'eau} (solide $<$ liquide $<$ gaz).
  \item L'état \rep{gazeux} a l'entropie la plus grande (molécules libres, très désordonnées) ; l'état
        \rep{solide} la plus petite (réseau ordonné).
\end{enumerate}
\end{correction}

\begin{correction}[signe de $\dS$ pour un changement d'état]
\begin{enumerate}[label=\alph*),leftmargin=2em,topsep=2pt,itemsep=2pt]
  \item Fusion (solide $\to$ liquide) : désordre $\nearrow$, \rep{$\dS>0$}.
  \item Vaporisation (liquide $\to$ gaz) : désordre $\nearrow$, \rep{$\dS>0$}.
  \item Condensation (gaz $\to$ liquide) : désordre $\searrow$, \rep{$\dS<0$}.
  \item Congélation (liquide $\to$ solide) : désordre $\searrow$, \rep{$\dS<0$}.
\end{enumerate}
\end{correction}

\begin{correction}[dissolution et précipitation]
\begin{enumerate}[label=\alph*),leftmargin=2em,topsep=2pt,itemsep=2pt]
  \item Dissolution : \rep{$\dS>0$}.
  \item Précipitation : \rep{$\dS<0$}.
  \item Dissoudre un solide ordonné disperse ses ions dans toute la solution : le désordre \emph{augmente}
        ($\dS>0$). La précipitation fait l'inverse : des ions dispersés se rangent dans un solide ordonné,
        le désordre \emph{diminue} ($\dS<0$).
\end{enumerate}
\end{correction}

\begin{correction}[les principes]
\begin{enumerate}[label=\alph*),leftmargin=2em,topsep=2pt,itemsep=2pt]
  \item Deuxième principe : un processus est spontané si l'entropie de l'univers augmente,
        \rep{$\dS_{\text{univers}}>0$}.
  \item À $T=\SI{0}{\kelvin}$, l'entropie d'un cristal parfait est \rep{$S=\SI{0}{\joule\per\mol\per\kelvin}$} :
        c'est le \rep{troisième principe}.
\end{enumerate}
\end{correction}

\begin{correction}[le désordre augmente-t-il ?]
\begin{enumerate}[label=\alph*),leftmargin=2em,topsep=2pt,itemsep=2pt]
  \item \rep{Vrai} : sublimation (solide $\to$ gaz), le désordre augmente nettement, $\dS>0$.
  \item \rep{Faux} : on passe de $2$ mol de gaz à $1$ mol de gaz ($2\,\ce{NO2}\to\ce{N2O4}$), le désordre
        \emph{diminue}, $\dS<0$.
  \item \rep{Vrai} : on crée du gaz ($\ce{CO2}$) à partir d'un solide, le désordre augmente, $\dS>0$.
\end{enumerate}
\end{correction}

\end{document}
